% Part:sets-functions-relations
% Chapter: sets
% Section: schroder-bernstein

\documentclass[../../../include/open-logic-section]{subfiles}

\begin{document}

\olfileid{sfr}{siz}{sb}
\olsection{Het begrip grootte en Schr\"oder-Bernstein}
\begin{explain}
Een intuïtieve gedachte is de volgende: als $A$ niet groter is dan $B$
en $B$ niet groter is dan $A$, dan zijn $A$ en $B$ gelijkmachtig.
Eerlijk gezegd zouden we, als deze gedachte \emph{onjuist} was,
moeilijk kunnen volhouden dat het door ons gedefinieerde begrip van
gelijkmachtigheid iets te maken heeft met het vergelijken van de
‘groottes’ van verzamelingen!{} Gelukkig klopt de intuïtieve gedachte.
Dit volgt uit de stelling van Schr\"oder-Bernstein.
\end{explain}
\begin{thm}[Schr\"oder-Bernstein]
	\ollabel{thm:schroder-bernstein}
	Als $\cardle{A}{B}$ en $\cardle{B}{A}$,
	dan $\cardeq{A}{B}$.
\end{thm}
\begin{explain}
Met andere woorden: als er !!a{injection} van $A$ naar~$B$ bestaat en
!!a{injection} van $B$ naar~$A$, dan bestaat er !!a{bijection} van $A$
naar~$B$.
Dit resultaat is echter tamelijk \emph{moeilijk} te bewijzen. Inderdaad:
Cantor formuleerde het resultaat, maar anderen hebben het
bewezen.\footnote{Zie voor meer over de geschiedenis bijvoorbeeld
\citet[pp.~165--6]{Potter2004}.}
\oliflabeldef{sfr:cardinals:card-sb:sec}{We kunnen pas
\emph{bewijzen} dat Schr\"oder-Bernstein geldt in
\olref[sfr][cardinals][card-sb]{sec}.}{}%
Voorlopig kunt (en moet) u dit zonder bewijs aannemen.
Gelukkig is Schr\"oder-Bernstein \emph{juist}. Daarmee bevestigt de
stelling onze gedachte dat de door ons gedefinieerde relaties, namelijk
$\cardeq{A}{B}$ en $\cardle{A}{B}$, iets met ‘grootte’ te maken hebben.
Bovendien is Schr\"oder-Bernstein zeer \emph{nuttig}. Het kan lastig
zijn !!a{bijection} tussen twee gelijkmachtige verzamelingen te
vinden. Met de stelling van Schr\"oder-Bernstein kunnen we die opgave in
twee richtingen opsplitsen: we hoeven alleen !!a{injection} van de eerste
naar de tweede verzameling te vinden en vervolgens hetzelfde in de
andere richting te doen.
\end{explain}
\end{document}